\section[曲线上向量丛的叠]{曲线上向量丛的叠（11 月 30 日）}

\srcpage{49}
\subsection{向量丛的 v-下降}

设 $E$ 为非阿基米德局部域，剩余域为 $\Fq$，$\pi\in\cO_E\subset E$ 是一致化子。
对 $S\in\Perf_{\Fq}$，有相对 Fargues--Fontaine 曲线 $X_S=X_{S,E}$。

\begin{definition}
固定 $n\ge 1$。令 $\Bun_n$ 为 $\Perf_{\Fq}$ 上的模（预）叠，它把
\[
S\longmapsto\{\text{$X_S$ 上秩为 $n$ 的向量丛}\}.
\]
\end{definition}

\begin{proposition}[{[Berk, Lemma 17.1.8]}]
\begin{enumerate}[label=(\arabic*),itemsep=1pt,topsep=2pt]
\item $\Bun_n$ 是 v-叠。
\item 在 $\mathrm{Perfd}:=\{\text{$\Zp$ 上的完满空间}\}$ 上，向量丛满足 v-下降。
\end{enumerate}
\end{proposition}

\begin{proof}
(2) 由 Kedlaya--Liu 已知解析下降。于是可化到 affinoid 的情形。要证的是：若
$Y=\Spa(S,S^+)\to X=\Spa(R,R^+)$ 是 affinoid 完满空间之间的 v-覆盖（这里指
$|Y|\twoheadrightarrow|X|$），记 $\Proj(R)$ 为有限投射 $R$-模的范畴，则函子
\[
\Proj(R)\xrightarrow{\ \sim\ }
\left\{N\in\Proj(S)\ \middle|\
\begin{array}{c}
\alpha\colon N\widehat\otimes_R S\xrightarrow{\sim}S\widehat\otimes_R N\
\text{是有限投射 }S\widehat\otimes_R S\text{-模的同构，}\\[2pt]
\text{且在 }S\widehat\otimes_R S\widehat\otimes_R S\text{ 上满足上闭链条件}
\end{array}\right\}
\]
是范畴的等价。

完全忠实性已知，因为有右伴随
\[
\operatorname{eq}(N\Rightarrow N\otimes_R S)\longleftarrow(N,\alpha).
\]
伴随的单位 $M\mapsto\operatorname{eq}(M\widehat\otimes_R S\Rightarrow
M\otimes_R S\widehat\otimes_R S)$ 是同构，因为结构层是 v-层（这给出 $M=R$
的情形，再与 $M$ 作张量即得一般情形）。

第一步：处理 $R=K$ 为完满域的情形。此时 $S$ 是 v-覆盖当且仅当 $S\ne 0$。
可设 $S$（一个 $K$-Banach 代数）是拓扑上可数生成的。这是因为任何 $K$-Banach
代数都是其拓扑上可数生成的子空间的滤过余极限，而完备化张量积与可数滤过余极限交换，
因为完备化过程只添加可数和。

这样做的好处是：可数生成的 Banach 空间作为 Banach 空间是自由的，故 $S$ 作为
$K$-Banach 空间自由。特别地，$(-)\widehat\otimes_K S$ 正合且保守，因为 $K$ 是
$S$ 的直和项。

设 $(N,\alpha)$ 是下降数据，$M=\operatorname{eq}(N\Rightarrow N\widehat\otimes_K S)$。要证
\[
M\widehat\otimes_K S\xrightarrow{\ \sim\ }N.
\]
但 $M\widehat\otimes_K S=\operatorname{eq}(N\widehat\otimes_K S\Rightarrow
N\widehat\otimes_K S\widehat\otimes_K S)$，而该等化子总是 $N$，因为
\[
0\to N\to N\widehat\otimes_K S\to N\widehat\otimes_K S\widehat\otimes_K S\to\cdots
\]
总是正合的：它有收缩同伦。（要点在于，这些完备化张量积是 $K$-Banach 空间范畴中的
张量积，于是这就是通常的形式论证。）

第二步：处理一般的 $R\to S$。取任意点 $x\in X=\Spa(R,R^+)$，其完备剩余域为
$K(x)$（某个完满域）。基变换给出
\[
K(x)\to S\widehat\otimes_R K(x),
\]
沿它可用第一步作下降。特别地，给定任意下降数据 $(N,\alpha)$，
$\bigl(N\widehat\otimes_S(S\widehat\otimes_R K(x))\bigr)$ 有限自由，且承认一组在
$\alpha$ 下不变的基。

\srcpage{50}
于是对 $x$ 的某个小有理邻域 $U\subset X$，
$N\widehat\otimes_S(S\widehat\otimes_R\cO_X(U))$ 有限自由，且承认一组基，使 $\alpha$
在该基下的矩阵 $\equiv 1\pmod{\varpi}$，其中 $\varpi\in R^+$ 是伪一致化子。
理由是：可以先在有理邻域上逼近 $x$ 处的一组基，再由标准的极限论证知它在某个小邻域上
仍是基；而「$\equiv1\pmod\varpi$」是开条件，在 $x$ 处成立，故可延拓到某个开邻域。

由解析下降可设 $U=X$。于是 $(N,\alpha)=(S^n,\alpha\in\mathrm{GL}_n(S\widehat\otimes_R S))$。
前面的化约保证了 $\alpha\in\mathrm{GL}_n(S^+\widehat\otimes_{R^+}S^+)$ 且
$\alpha\equiv1\pmod\varpi$。要找一个基变换使 $\alpha$ 成为 $\mathrm{Id}$。

我们用逐次逼近来证明该断言。注意
\[
\frac{\alpha-1}{\varpi}\pmod{\varpi}\in M_n(S^+\widehat\otimes_{R^+}S^+/\varpi)
\]
是加性上闭链。但 $H^1_v(X,\cO^+/\varpi)$ 几乎为 $0$，即被 $\varpi^\epsilon$ 杀死
（任意 $\epsilon>0$）。因此可作基变换使
$\alpha\equiv1\pmod{\varpi^{2-\epsilon}}$（任意 $\epsilon>0$），再归纳地继续。

(1) 要证 $\Bun_n$ 是 v-叠。这里用到：$X_S\times_{\Spa E}\Spa E_\infty$ 是完满的，
其中 $E_\infty=E(\pi^{1/p^\infty})^\wedge$，且这一构造把 v-覆盖变为 v-覆盖。
于是可用 (2) 下降 $X_S\times_{\Spa E}\Spa E_\infty$ 上的向量丛；再按情形 1 的论证
沿 $E_\infty/E$ 下降。
\end{proof}

问题：能否也下降完美复形？我们预期答案是肯定的。

\begin{remark}
在 [BS17] 中，对完美概形上的向量丛与完美复形也证明了类似的 v-下降结果。
\end{remark}

\subsection{$\Bun_n$ 的结构}

由向量丛的分类，$|\Bun_n|$ 只有可数多个点，它们由宽度为 $n$、拐点坐标为整数的
Newton 多边形枚举。

\begin{center}
\figP{050}
\end{center}

记 $B(\mathrm{GL}_n)$ 为这类 Newton 多边形的集合，亦即
$\{n\text{ 维等晶体}\}/\sim$。于是得到双射
$|\Bun_n|=B(\mathrm{GL}_n)$。

$|\Bun_n|$ 上有自然的拓扑：称 $U\subset|\Bun_n|$ 为开，若它对应一个开子叠。
等价地说，若 $X\twoheadrightarrow\Bun_n$ 是完满空间给出的 v-覆盖，且
$Y\twoheadrightarrow X\times_{\Bun_n}X$，则 $|\Bun_n|=|X|/|Y|$。这个拓扑的结构
是什么？

在 $B(\mathrm{GL}_n)$ 上引入偏序，即支配序：若 $P$ 位于 $P'$ 之上或与之重合，且端点
相同，则称 $P\ge P'$。这诱导 $B(\mathrm{GL}_n)$ 上的一个拓扑：$U\subset B(\mathrm{GL}_n)$
为开，若对一切 $P\in U$，所有 $P'\ge P$ 也落在 $U$ 中。

\srcpage{51}
定理 9.4 说明 $|\Bun_n|\to B(\mathrm{GL}_n)$ 连续。

\begin{theorem}[{[BFHHLWY]}]
映射 $|\Bun_n|\to B(\mathrm{GL}_n)$ 是同胚。
\end{theorem}

\begin{remark}
对一般的 $G$，Viehmann 最近用完全不同的方法宣布了该结果。[BFHHLWY] 的方法此前已被
Hamann 推广到一些典型群。
\end{remark}

\begin{example}
$|\Bun_1|=|\Pic|\cong\ZZ$ 离散，其点为 $\cO,\cO(1),\dots$
\begin{center}
\figP{051A}
\end{center}
\end{example}

\begin{example}
来看 $|\Bun_2|$。端点给出一个连通分支，在这里它度量次数 $d$。对每个 $d$，存在一个
半稳定丛。
\begin{center}
\figP{051B}
\end{center}
一般点的余维为 $0$，而特化的余维越来越高。与 $\cO(1)$ 作张量使图形在竖直方向以 $2$
为周期。
\end{example}

\srcpage{52}
\begin{example}
我们画出 $|\Bun_3|$ 的 $d=0$ 分层的开端。
\begin{center}
\figP{052}
\end{center}
\end{example}

\begin{remark}
这看起来有点像 $\PP^1$ 上的 $\Bun_n$。
\end{remark}

观察：$\pi_0(\Bun_n)\xrightarrow{\sim}\ZZ$ 由次数给出；每个连通分支有唯一的半稳定点
$\cO(d/n)=\cO(a/b)^{\oplus c}$，其中 $d/n=ac/bc$ 且 $(a,b)=1$。

\subsection{分层的描述}

对 $b\in B(\mathrm{GL}_n)$，记 $\Bun^b_G\subset\Bun_n$ 为相应的局部闭分层。

\subsubsection{平凡分层}

先设 $b\cong\cO^n$。由定理 9.4 可得：对 $S\in\Perf_{\Fq}$，$X_S$ 上在每个纤维中
斜率为 $0$ 的半稳定向量丛，通过
\[
\cE\mapsto BC(\cE),\qquad \mathbb{L}\otimes_E\cO_{X_S}\longleftarrow\mathbb{L}
\]
与 $E$-局部系统等价。这来自 pro-étale 下降，以及如下事实：任何这样的 $\cE$ 在
pro-étale 局部平凡。

\begin{corollary}
对应于 $\cO^n$ 的 $\Bun_n$ 的分层是 $[\mathrm{pt}/\mathrm{GL}_n(E)]$，即分类
$E$-局部系统的叠。
\end{corollary}

（对任何拓扑空间 $T$，函子 $T\colon S\mapsto\operatorname{Cont}(|S|,T)$ 是 v-层。）

于是有开浸入 $j\colon[*/\mathrm{GL}_n(E)]\hookrightarrow\Bun_n$。特别地，
$\mathrm{GL}_n(E)$ 的表示是 $[*/\mathrm{GL}_n(E)]$ 上的层，可用 $j_!$ 将其嵌入
$\Bun_n$ 上的层范畴。

\subsubsection{半稳定丛}

接着分析半稳定点 $\cO(d/n):=\cO(a/b)^{\oplus c}$。此时
$\ul{\Aut}(\cO(d/n))=\mathrm{GL}_c(D_{a/b})$，其中 $D_{a/b}=\End(\cO(a/b))$ 是
$E$ 上 Hasse 不变量为 $a/b$ 的中心可除代数。

于是 $\Bun^b_n\cong[*/\mathrm{GL}_c(D_{a/b})]$。注意 $\mathrm{GL}_c(D_{a/b})$ 是
$\mathrm{GL}_n/E$ 的某个内形式的 $E$-值点。$\mathrm{GL}_n$ 的所有内形式都这样得到。
因此 $j_b\colon B(\mathrm{GL}_c(D_{a/b}))\to\Bun_n$ 诱导了
$\Rep(\mathrm{GL}_c(D_{a/b}))$ 到 $\Bun_n$ 上的层中的嵌入。所以 $\Bun_n$ 能看到
$\mathrm{GL}_n(E)$ 及其所有内形式的表示。

\begin{remark}
若要得到这些投射有限稳定化子，就需要在 pro-étale 拓扑下工作。事实上，
$\ast\to\ast/G$ 满射的含义是 $G$-主丛可在相应拓扑下平凡化；若 $G$ 是投射有限的，
则这只能在投射有限拓扑下做到。
\end{remark}

\srcpage{53}
\subsubsection{非半稳定点}

总有 $\Bun^b_n=[\mathrm{pt}/\ul{\Aut}(\cE_b)]$，其中 $\cE_b$ 是 $b$ 对应的向量丛，
$\ul{\Aut}(\cE_b)$ 是 v-层 $S\mapsto\Aut(\cE_b|_{X_S})$。

\begin{example}
来看 $b\cong\cO\oplus\cO(1)$ 时的 $\Bun_2$。此时
\[
\ul{\Aut}(\cO\oplus\cO(1))=
\begin{pmatrix}
E^\times & BC(\cO(1))\\
0 & E^\times
\end{pmatrix}
\]
特别地，回忆 $BC(\cO(1))$ 是完满开单位圆盘。特别地，$\ul{\Aut}(\cO\oplus\cO(1))$ 是
正维数的。它带有自然的滤过
\[
1\to\underbrace{BC(\cO(1))}_{1\text{ 维、连通}}\to
\ul{\Aut}(\cO\oplus\cO(1))\to
\underbrace{E^\times\times E^\times}_{0\text{ 维}}\to1.
\]
一般地，写 $\cE_b=\bigoplus_{\lambda\in\QQ}\cO(\lambda)^{n_\lambda}$，其中斜率
$\lambda$ 按递增顺序排列，并记 $\cE_b^\lambda$ 为相应的斜率 $\lambda$ 的部分。则
\[
\ul{\Aut}(\cE_b)=
\begin{pmatrix}
\ul{\Aut}(\cE_b^{\lambda_1}) & BC(\Hom(\cE_b^{\lambda_1},\cE_b^{\lambda_2})) & \\
 & \ul{\Aut}(\cE_b^{\lambda_2}) & BC(\Hom(\cE_b^{\lambda_2},\cE_b^{\lambda_3})) \\
 & & \ddots
\end{pmatrix}
\]
于是一般地有扩张
\[
1\to\underbrace{\text{「幂幺的」}}_{\text{维数 }>0\text{ 的 Banach--Colmez 空间之扩张}}
\to\ul{\Aut}(\cE_b)\to
\underbrace{\text{局部 profinite 群}}_{\ul{\Aut}(V_b)}\to1.
\]
其中 $V_b$ 是 $b$ 对应的等晶体。传统上记 $J_b(E)=\Aut(V_b)$。它是 $\mathrm{GL}_n$
的某个 Levi 子群的内形式的 $E$-值点。
\end{example}

\begin{remark}
幂幺群不能作用在 $\ell$-进层上，因此
\[
\{\Bun^b_n\text{ 上的 }\ell\text{-进层}\}\cong\{J_b(E)\text{ 的 }\ell\text{-进表示}\}.
\]
从而 $\{\Bun_n\text{ 上的 }\ell\text{-进层}\}$ 由形如
$\{J_b(E)\text{ 的 }\ell\text{-进表示}\}$ 的部件经半正交分解拼成。
\end{remark}

\subsection{分层之间如何相互作用？}

\begin{example}[$n=2$]
可以如下构造映射 $(\PP^1_E)^\diam\to\Bun_2$。一个映射
$S\to(\PP^1_E)^\diam$ 等价于一个 untilt $S^\#/E$ 加上一个满射
$\cO^2_{S^\#}\twoheadrightarrow L$。由这些数据可以构造向量丛
\[
\cE(L):=\ker\Bigl(\cO^2_{X_S}\to
\underbrace{\cO^2_{S^\#}}_{\text{来自 }S^\#\hookrightarrow X_S}
\twoheadrightarrow L\Bigr).
\]
它给出 $X_S$ 上秩 $2$ 的向量丛，从而是映射 $S\to\Bun_2$。
\end{example}

\begin{proposition}
像落在 $\Bun_2^{\cO(-1/2)}\cup\Bun_2^{\cO\oplus\cO(-1)}$ 中，其中
\begin{itemize}[itemsep=1pt,topsep=2pt]
\item $(\Omega^2)^\diam:=\PP^1_E\setminus\PP^1(E)$ 落在 $\Bun_2^{\cO(-1/2)}$ 中；
\item $\PP^1(E)$ 落在 $\Bun_2^{\cO\oplus\cO(-1)}$ 中。
\end{itemize}
\end{proposition}

\begin{proof}
\srcpage{54}对几何点 $S$，$\cE(L)$ 必为秩 $2$、次数 $-1$。于是
$\cE(L)\approx\cO(-1/2)$，或 $\cO(-i-1)\oplus\cO(i)$（某个 $i\ge0$）。但又有
$\cE(L)\hookrightarrow\cO^2$，故只可能 $i=0$。

若 $\cE(L)\cong\cO(-1)\oplus\cO$，则得到
\[
\cO\hookrightarrow\cE(L)\hookrightarrow\cO^2\twoheadrightarrow
\cO^2_{S^\#}\twoheadrightarrow L,
\]
取整体截面即得非零映射 $E\to E^2\to C^2\twoheadrightarrow L$。于是存在一条有理直线
被到 $L$ 的商杀死，故 $L$ 必为 $E$-有理的，从而给出 $\PP^1(E)$ 的一个点。反之，若该点
落在 $\PP^1(E)$ 中，则经 $\mathrm{GL}_2(E)$-作用可设其为 $[0:1]$，此时直接看出
$\cE(L)\cong\cO\oplus\cO(-1)$。
\end{proof}

\begin{remark}
映射 $(\PP^1_E)^\diam\to\Bun_2$ 接近于「光滑」。实际成立的是：
$(\PP^1_E)^\diam/\mathrm{GL}_2(E)\to\Bun_2$ 光滑，从而给出一个图卡。
\end{remark}
